\chapter{Propositions and Connectives}
\label{ch:propositions}

This chapter introduces the basic propositional connectives
$\land$, $\lor$, $\lnot$, $\to$, and $\leftrightarrow$, together with
Lean's \texttt{Bool}-valued truth-table decidability through the
\texttt{decide} tactic. It establishes, in miniature, the pairing
that organizes the whole book: a claim is proved where it holds, and
a plausible-looking converse or strengthening is shown false by an
explicit countermodel where it does not.

\section{The Basic Connectives}

Lean has two different homes for the five connectives, and keeping
them straight matters more than it might first appear.

At the level of \texttt{Prop}, each connective is a type whose
inhabitants are proofs, not a computed value: $p \land q$ is the
structure \texttt{And}, whose single constructor
\texttt{And.intro} packages a proof of $p$ together with a proof of
$q$; $p \lor q$ is the inductive type \texttt{Or}, with two
constructors \texttt{Or.inl} and \texttt{Or.inr} for proving it from
either side (already used in the worked example below); $\lnot p$ is
literally notation for $p \to \mathrm{False}$, so a proof of
$\lnot p$ is a function turning any proof of $p$ into a proof of
\texttt{False}; $p \to q$ needs no special type at all, since a proof
of an implication simply \emph{is} a function from proofs of $p$ to
proofs of $q$; and $p \leftrightarrow q$ is the structure
\texttt{Iff}, bundling a proof of $p \to q$ with a proof of
$q \to p$.

At the level of \texttt{Bool}, the same five symbols name ordinary
total computable functions on the two-element type
$\{\mathtt{true}, \mathtt{false}\}$: \texttt{\&\&}, \texttt{||},
\texttt{!}, and Boolean implication all just look up an answer in a
finite table. Because that table is finite, equality between two
\texttt{Bool}-valued expressions is always decidable — which is
exactly why \texttt{decide} can close the countermodel below by
brute enumeration, with nothing constructive required.
\begin{lean}
-- Prop level: a proof is data, not a computed answer.
example (p q : Prop) (hp : p) (hq : q) : p ∧ q := And.intro hp hq

-- Bool level: a value is computed by table lookup.
example : (true && false) = false := rfl
\end{lean}
The \texttt{decide} tactic is the bridge between the two registers:
for a \texttt{Prop} that happens to be decidable, it computes the
\texttt{Bool} answer and uses that computation directly as the
proof. Every \texttt{decide} call in this book is doing exactly
that, silently, underneath the enumeration it appears to perform.

\section{A Worked Pair: Disjunction Introduction}

The first genuine result of the book is small on purpose: it should
be possible to state, prove, and understand the countermodel to its
converse within a single sitting.

\begin{experiment}{Disjunction introduction}
For any propositions $p$ and $q$, $p \to p \lor q$.
\begin{lean}
theorem or_intro_left (p q : Prop) (hp : p) : p ∨ q :=
  Or.inl hp
\end{lean}
The proof is definitional: \texttt{Or.inl} is literally the
introduction rule for the left disjunct. Nothing about $q$ is used
or needed.
\end{experiment}

\begin{countermodel}{The converse of disjunction introduction}
The converse, $p \lor q \to p$, is not valid: it fails whenever $q$
holds and $p$ does not.
\begin{lean}
example : ¬ (∀ p q : Bool, (p || q) = true → p = true) := by
  decide
\end{lean}
Working at the level of \texttt{Bool} rather than \texttt{Prop}
makes the countermodel a finite, decidable check: Lean simply
enumerates the four truth-value combinations and exhibits
$p = \mathrm{false}, q = \mathrm{true}$ as a witness where the
premise holds and the conclusion fails.
\end{countermodel}

\section{De Morgan, Contraposition, Distributivity}

Contraposition itself — $(p \to q) \to (\lnot q \to \lnot p)$ — is
left as Exercise~2 below rather than worked here, since attempting
it after this section's examples is a reasonable checkpoint. What
follows instead are De Morgan's laws and distributivity, plus one
plausible-looking strengthening of distributivity that is not valid.

\begin{experiment}{De Morgan for $\lor$, both directions}
Unlike some non-classical distinctions the book will return to in
Chapter~\ref{ch:classical-constructive}, this particular law needs
no classical principle in either direction.
\begin{lean}
theorem demorgan_or (p q : Prop) : ¬(p ∨ q) ↔ (¬p ∧ ¬q) :=
  Iff.intro
    (fun h => And.intro
      (fun hp => h (Or.inl hp)) (fun hq => h (Or.inr hq)))
    (fun h => fun hpq => hpq.elim h.1 h.2)
\end{lean}
Both halves of the \texttt{Iff} are ordinary functions built from
the constructors already in hand — written here with
\texttt{Iff.intro}/\texttt{And.intro} spelled out rather than the
usual $\langle \cdot, \cdot \rangle$ shorthand, since nesting that
shorthand two levels deep does not render reliably in this book's
code-listing setup: the forward direction turns a disproof of the
whole disjunction into two separate disproofs by composing with
\texttt{Or.inl}/\texttt{Or.inr}; the backward direction uses
\texttt{Or.elim} (written here via dot notation, \texttt{hpq.elim})
to case on which disjunct held.
\end{experiment}

\begin{experiment}{De Morgan for $\land$, one direction}
The direction from a disjunction of negations to a negated
conjunction needs nothing beyond what is already available.
\begin{lean}
theorem demorgan_and_one_way (p q : Prop) (h : ¬p ∨ ¬q) :
    ¬(p ∧ q) :=
  fun hpq => h.elim (fun hnp => hnp hpq.1) (fun hnq => hnq hpq.2)
\end{lean}
The converse — from $\lnot(p \land q)$ to $\lnot p \lor \lnot q$ —
is also true, but not by a proof this direct: it needs a classical
principle unavailable at this point in the book, and is revisited
once Chapter~\ref{ch:classical-constructive} has introduced one.
\end{experiment}

\begin{experiment}{Distributivity of $\land$ over $\lor$}
\begin{lean}
theorem and_distrib_or (p q r : Prop) :
    p ∧ (q ∨ r) ↔ (p ∧ q) ∨ (p ∧ r) :=
  Iff.intro
    (fun h => h.2.elim
      (fun hq => Or.inl (And.intro h.1 hq))
      (fun hr => Or.inr (And.intro h.1 hr)))
    (fun h => h.elim
      (fun hpq => And.intro hpq.1 (Or.inl hpq.2))
      (fun hpr => And.intro hpr.1 (Or.inr hpr.2)))
\end{lean}
\end{experiment}

\begin{countermodel}{Distributivity does not extend to full separation}
It is tempting to guess that $p \land (q \lor r)$ forces both
$p \land q$ \emph{and} $p \land r$ separately, not merely one or the
other. It does not.
\begin{lean}
example : ¬ (∀ p q r : Bool,
    (p && (q || r)) = true → ((p && q) && (p && r)) = true) := by
  decide
\end{lean}
Taking $p = \mathtt{true}$, $q = \mathtt{true}$, $r = \mathtt{false}$
satisfies the premise ($p$ holds and at least one of $q, r$ does) but
fails the conclusion ($p \land r$ is false), exactly the witness
\texttt{decide} finds by exhaustive search over the eight cases.
\end{countermodel}

\section{Application: Digital Logic Circuits}

Propositional logic is not only a tool for evaluating arguments; it
is, quite literally, the specification language for digital
hardware. An AND gate, an OR gate, and a NOT gate compute $\land$,
$\lor$, and $\lnot$ on electrical signals instead of truth values,
and a circuit built from them computes whatever Boolean expression
its wiring diagram represents. This gives the truth tables and
\texttt{decide} calls from earlier in the chapter a second life: two
circuits that look completely different can be shown to compute the
exact same function by checking that their Boolean expressions agree
on every input — which is exactly what \texttt{decide} already does.

\begin{experiment}{Reading a circuit off a truth table}
Given a truth table, a Boolean expression that matches it can always
be built directly: for each row where the output is true, form the
conjunction of each input (negated if that input is false in that
row), then join all such conjunctions with $\lor$. For the table
\[
\begin{array}{ccc|c}
P & Q & R & S \\ \hline
1 & 1 & 1 & 1 \\
1 & 1 & 0 & 0 \\
1 & 0 & 1 & 1 \\
1 & 0 & 0 & 1 \\
0 & 1 & 1 & 0 \\
0 & 1 & 0 & 0 \\
0 & 0 & 1 & 0 \\
0 & 0 & 0 & 0
\end{array}
\]
the three true rows give
$S = (P \land Q \land R) \lor (P \land \lnot Q \land R) \lor (P \land \lnot Q \land \lnot R)$.
Rather than trust that construction by eye, state the table directly
as one function and the built expression as another, then let
\texttt{decide} confirm they agree everywhere:
\begin{lean}
def sFromTable (p q r : Bool) : Bool :=
  match p, q, r with
  | true,  true,  true  => true
  | true,  true,  false => false
  | true,  false, true  => true
  | true,  false, false => true
  | false, true,  true  => false
  | false, true,  false => false
  | false, false, true  => false
  | false, false, false => false

def sFromDNF (p q r : Bool) : Bool :=
  (p && q && r) || (p && !q && r) || (p && !q && !r)

example : ∀ p q r : Bool, sFromTable p q r = sFromDNF p q r := by
  decide
\end{lean}
\end{experiment}

\begin{experiment}{Two circuits, one function}
Circuits with visibly different wiring can compute identical
functions — the point of a later stage in circuit design is finding
the simpler one. A circuit built from an OR of two ANDs need be no
more powerful than a single wire:
\begin{lean}
def circuitA (p q : Bool) : Bool := (p && q) || (p && !q)
def circuitB (p q : Bool) : Bool := p

example : ∀ p q : Bool, circuitA p q = circuitB p q := by decide
\end{lean}
\texttt{circuitA} and \texttt{circuitB} disagree on how many gates
they use and agree on every output — the same distinction between
syntactic form and semantic content that separates a Lean proof
term from the proposition it proves.
\end{experiment}

\section{Exercises}

\begin{exercise}
Prove $q \to p \lor q$ (disjunction introduction, right).
\end{exercise}

\begin{exercise}
State and prove, or refute by countermodel, the claim
$(p \to q) \to (\lnot q \to \lnot p)$.
\end{exercise}

\begin{exercise}
Using \texttt{decide} on \texttt{Bool}, determine whether
$(p \land q) \lor r$ and $(p \lor r) \land (q \lor r)$ always agree.
\end{exercise}
